% BHU PYQ -- Manually transcribed & error-corrected
% Paper: BCH-111 : Financial Accounting-I
% Exam: B.Com. (Hons.) Semester I Examination, 2023-24

\documentclass[11pt,a4paper]{article}
\usepackage[a4paper,top=2.5cm,bottom=2.5cm,left=2.8cm,right=2.8cm,headheight=14pt]{geometry}
\usepackage{amsmath,amssymb}
\usepackage[T1]{fontenc}
\usepackage[utf8]{inputenc}
\usepackage{lmodern,microtype,fancyhdr,enumitem,hyperref,lastpage,parskip}
\usepackage{booktabs}

\hypersetup{
    colorlinks=true,
    urlcolor=black,
    linkcolor=black,
    pdftitle={BCH-111: Financial Accounting-I -- BHU B.Com. (Hons.) Sem I 2023-24}
}

\pagestyle{fancy}
\fancyhf{}
\renewcommand{\headrulewidth}{0.4pt}
\renewcommand{\footrulewidth}{0.4pt}
\fancyhead[L]{\small   \textit{Banaras Hindu University}}
\fancyhead[R]{\small   \textit{BCH-111: Financial Accounting-I}}
\fancyfoot[C]{\small Page  \thepage\ of \pageref{LastPage}}


\newlist{parts}{enumerate}{2}
\setlist[parts,1]{label=(\alph*),leftmargin=*,itemsep=5pt,topsep=3pt}
\setlist[parts,2]{label=(\roman*),leftmargin=*,itemsep=3pt,topsep=2pt}

\newcommand{\pts}[1]{\hfill   \textit{[#1]}}


\begin{document}


\begin{center}
  {\large\bfseries BANARAS HINDU UNIVERSITY}\\[2pt]
  {
\normalsize B.Com. (Hons.) --- Semester I Examination, 2023-24}\\[6pt]
  \rule{0.8\linewidth}{0.5pt}\\[6pt]
  {\large\bfseries Paper No. BCH-111: Financial Accounting-I}\\[4pt]
  {
\normalsize Time: 3 Hours\hfill Maximum Marks: 70}
\end{center}

\rule{\linewidth}{0.4pt}

\smallskip



\noindent   \textit{Instructions: Answer all questions. All questions carry equal marks (14 marks each).}

\bigskip




\noindent   \textbf{Question 1.}
Journalise the following transactions in the books of M/s ABC for the month of April 2023:\pts{14}

\begin{itemize}[itemsep=2pt,topsep=2pt]
  \item[April 1:] Opening balances in different accounts: Leasehold Premises \pounds 3,80,000; Machinery \pounds 2,76,000; Stock \pounds 5,19,000; Cash at Bank \pounds 76,200; Cash in Hand \pounds 8,600; Loan @ 12\% \pounds 2,00,000; Creditors \pounds 2,00,000; Due to Wilson \& Grey \pounds 60,000; Due from Mohan \& Co. \pounds 14,600; Due from Cavendish \& Co. \pounds 12,600; Furniture \pounds 15,000.
  \item[April 1:] Infused additional capital in business: \pounds 5,00,000.
  \item[April 5:] Received a cheque from Cavendish \& Co. in full settlement of his balance till date: \pounds 12,350.
  \item[April 10:] Sold goods to Cavendish \& Co.: \pounds 30,400.
  \item[April 18:] Purchased a motorcycle for the owner: \pounds 50,000.
  \item[April 25:] Cavendish \& Co. went insolvent. Nothing is recoverable from them.
  \item[April 30:] Allow interest (rounded to the nearest rupee) on capital @ 10\% per annum.
\end{itemize}

\centerline{   \textbf{OR}}

Write notes on any two of the following:\pts{14}

\begin{parts}
  \item Capital Expenditure and Revenue Expenditure
  \item Subsidiary Books
  \item Dual Aspect Concept
\end{parts}

\medskip\rule{\linewidth}{0.2pt}


\newpage




\noindent   \textbf{Question 2.}
Following is the Trial Balance of Alfa Traders as on 31st March, 2021:


\begin{center}

\begin{tabular}{lr|lr}
 oprule
   \textbf{Debit Balances} &   \textbf{Amount (\pounds)} &   \textbf{Credit Balances} &   \textbf{Amount (\pounds)} \\
\midrule
Opening Stock & 1,00,000 & Capital & 1,60,000 \\
Sundry Debtors & 1,55,000 & Sundry Creditors & 50,000 \\
Cash at Bank & 15,000 & Sales & 4,50,000 \\
Sundry Assets & 2,93,000 & Discount and Commission & 12,500 \\
Bad Debts & 1,000 & Provision for Doubtful Debts & 10,000 \\
Accrued Commission & 1,700 & Salaries Outstanding & 12,000 \\
Purchases & 1,50,500 & Loan (interest rate 12\% p.a.) & 1,00,000 \\
Depreciation & 4,500 & Suspense Account & 1,200 \\
Income Tax (Drawings) & 10,000 & & \\
Salaries & 36,000 & & \\
Wages & 28,000 & & \\
Interest on Loan & 1,000 & & \\
\midrule
   \textbf{Total} &   \textbf{7,95,700} &   \textbf{Total} &   \textbf{7,95,700} \\
ottomrule
\end{tabular}
\end{center}




\noindent   \textbf{Adjustments:}

\begin{parts}
  \item Closing Stock at the end of the year was valued at \pounds 80,000.
  \item Further bad debts occurred during the year were \pounds 3,000.
  \item Maintain a 5\% provision for doubtful debts.
  \item An amount of \pounds 1,200 paid as rent was wrongly debited to the landlord's account and included in debtors.
  \item Dishonour of a cheque for \pounds 5,000 was not recorded in the books. It is expected that the concerned debtor will not pay more than 40\% of the amount due.
  \item The Suspense Account in the Trial Balance was on account of a wrong totaling of the discount received column in the cash book.
\end{parts}
Prepare the Trading and Profit \& Loss Account for the year ended 31st March, 2021 and a Balance Sheet as on that date.\pts{14}

\centerline{   \textbf{OR}}

What is a Trial Balance? Why and how is it prepared?\pts{14}

\medskip\rule{\linewidth}{0.2pt}


\newpage




\noindent   \textbf{Question 3.}
A and B are partners sharing profit in the ratio of 3:2. The Balance Sheet of their firm as at 31st March, 2022 is as follows:


\begin{center}

\begin{tabular}{lr|lr}
 oprule
   \textbf{Liabilities} &   \textbf{Amount (\pounds)} &   \textbf{Assets} &   \textbf{Amount (\pounds)} \\
\midrule
Creditors & 1,40,000 & Cash in Hand & 1,10,000 \\
Bank Overdraft & 2,00,000 & Sundry Debtors & 90,000 \\
Capital Accounts: & & Land and Building & 4,00,000 \\
~~A & 3,50,000 & Plant \& Machinery & 3,40,000 \\
~~B & 2,50,000 & & \\
\midrule
   \textbf{Total} &   \textbf{9,40,000} &   \textbf{Total} &   \textbf{9,40,000} \\
ottomrule
\end{tabular}
\end{center}




\noindent C is admitted as a partner with effect from 1st April, 2022 with the new profit sharing ratio being 2:2:1. The other details relating to C's admission are:

\begin{parts}
  \item C will bring in \pounds 2,00,000 as his share of capital.
  \item The value of the firm's goodwill is valued at \pounds 2,00,000.
  \item An amount of \pounds 5,000 owing to D for purchase of goods has been omitted from the list of Sundry Creditors.
  \item Building is to be revalued at \pounds 5,00,000 and Plant is to be written down to \pounds 3,00,000.
  \item The partners shall either pay cash to the firm or withdraw cash from the firm to adjust their capital balances in accordance with the new profit sharing ratio.
\end{parts}
Pass the necessary Journal Entries to record the admission of the new partner and prepare the Balance Sheet of the new firm.\pts{14}

\centerline{   \textbf{OR}}

What is a Memorandum Revaluation Account? Why and how is it prepared?\pts{14}

\medskip\rule{\linewidth}{0.2pt}


\newpage




\noindent   \textbf{Question 4.}
X, Y, and Z are partners in a firm. They share profits and losses in their capital ratio as it appeared in the Balance Sheet as on 31st March, 2021. On 1st April, 2021 Y retired. The Balance Sheet as on 31st March, 2021 is as under:


\begin{center}

\begin{tabular}{lr|lr}
 oprule
   \textbf{Liabilities} &   \textbf{Amount (\pounds)} &   \textbf{Assets} &   \textbf{Amount (\pounds)} \\
\midrule
Sundry Creditors & 2,30,000 & Cash at Bank & 50,000 \\
Capital Accounts: & & Stock & 1,80,000 \\
~~X & 6,00,000 & Sundry Debtors & 3,40,000 \\
~~Y & 4,00,000 & Land and Building & 5,10,000 \\
~~Z & 2,00,000 & Machinery & 3,50,000 \\
\midrule
   \textbf{Total} &   \textbf{14,30,000} &   \textbf{Total} &   \textbf{14,30,000} \\
ottomrule
\end{tabular}
\end{center}




\noindent The retirement of Y is based on the following terms:

\begin{parts}
  \item The new profit sharing ratio between X and Z shall be 3:2.
  \item Y's share of goodwill is to be adjusted through the remaining partners' capital accounts. The value of goodwill of the firm is \pounds 3,00,000.
  \item The capital amount payable to Y is to be paid in three equal annual installments after transferring his claim to his loan account. His loan account bears an interest rate of 6\% p.a.
  \item The value of Land and Building is to be appreciated by \pounds 3,30,000.
  \item An unrecorded contingent liability of \pounds 30,000 has devolved and is payable.
  \item Finally, the total capital balance of X and Z will be adjusted in their new ratio of 3:2. For this, X and Z will bring in or withdraw cash as the case may be.
\end{parts}
Prepare the necessary ledger accounts (Revaluation A/c, Partners' Capital A/cs, Y's Loan A/c) and the Balance Sheet of the reconstructed firm.\pts{14}

\centerline{   \textbf{OR}}

Explain the concept of goodwill and give journal entries for recording goodwill in each of the following cases:\pts{14}

\begin{parts}
  \item When a partner retires and a goodwill account is to be raised in the books.
  \item When a partner retires and a goodwill account is raised and subsequently written off.
  \item When the retiring partner's share is to be adjusted into the accounts of continuing partners without raising a Goodwill account.
\end{parts}

\medskip\rule{\linewidth}{0.2pt}


\newpage




\noindent   \textbf{Question 5.}
The Balance Sheet of a firm, when a partner D became insolvent on 31st March, 2021, stood as follows:


\begin{center}

\begin{tabular}{lr|lr}
 oprule
   \textbf{Liabilities} &   \textbf{Amount (\pounds)} &   \textbf{Assets} &   \textbf{Amount (\pounds)} \\
\midrule
Creditors & 1,50,000 & Sundry Fixed Assets & 2,50,000 \\
General Reserve & 1,20,000 & Cash in Hand and Bank & 40,000 \\
Capital Accounts: & & Capital Accounts: & \\
~~A & 1,30,000 & ~~C (debit balance) & 1,10,000 \\
~~B & 1,10,000 & ~~D (debit balance) & 1,10,000 \\
\midrule
   \textbf{Total} &   \textbf{5,10,000} &   \textbf{Total} &   \textbf{5,10,000} \\
ottomrule
\end{tabular}
\end{center}




\noindent Partners shared profits and losses equally. The fixed assets of the firm realized \pounds 2,00,000 and creditors were paid in full. You are required to close the books of the firm applying the Garner vs. Murray rule.\pts{14}

\centerline{   \textbf{OR}}

Give the accounting entries to record the dissolution of a partnership firm. Also explain the important points and applicability of the Garner vs. Murray rule.\pts{14}

\vfill
\rule{\linewidth}{0.4pt}

\begin{center}\small
\href{https://bhu-exam.vercel.app/}{Click here to visit our website}\\[4pt]
\href{https://me-aryan.vercel.app/}{Click here to visit our contributor}\\[4pt]
\href{https://quashan-me.vercel.app/}{Click here to visit our contributor}
\end{center}

\end{document}
